\documentclass[10pt,a4paper]{amsart}
\usepackage[margin=19mm]{geometry}
\usepackage{amsmath,amssymb,mathtools}
\usepackage[scr=rsfs,cal=euler]{mathalfa}
\usepackage{xfrac}
\usepackage[unicode,hidelinks,bookmarks=false]{hyperref}
\newcommand{\ie}{i.e.\ }
\newcommand{\cf}{cf.\ }
\newcommand{\resp}{respectively\ }
\newcommand{\Subsection}{\subsection}
\providecommand{\nobreakdash}{\nobreak\hbox{-}}
\setlength{\emergencystretch}{2em}
\multlinegap0pt
\usepackage{mathtools}
\usepackage{amssymb}
\usepackage{tikz}
\usepackage{xcolor}
\usepackage[shortlabels]{enumitem}
\newtheorem{theorem}{Theorem}
\newtheorem{lemma}{Lemma}
\newcommand{\CC}{\mathbb{C}}
\newcommand{\Mult}{\operatorname{Mult}}
\newcommand{\opphi}[4]{{#1}\in \Mult({#3},{#4})}
\title{\large BEURLING CRITERIA FOR REPRODUCING KERNELS}
\author{SHUAIBING LUO, SCOTT MCCULLOUGH, GEORGIOS TSIKALAS}
\date{Source: arXiv:2607.25088v2 (CC BY 4.0)}
\begin{document}
\maketitle

\noindent\emph{Source and licence.} \large BEURLING CRITERIA FOR REPRODUCING KERNELS. Version arXiv:2607.25088v2 (CC BY 4.0), distributed by arXiv under the Creative Commons Attribution 4.0 International licence, http://creativecommons.org/licenses/by/4.0/, as recorded in the arXiv licence metadata for this exact version and shown on the abstract page https://arxiv.org/abs/2607.25088v2. Full licence notice: \texttt{licenses/arxiv-2607.25088-CC-BY-4.0.txt}.\par\smallskip

\noindent\textit{Editorial scope.} Counted unit: the article's theorem realization. (source0.tex lines 1663--1675). The article's complete original proof of it is delivered with it (source0.tex lines 1678--1730, one \texttt{proof} environment of 53 lines). No mathematics is written, repaired or re-derived: the statement and the proof are the article's own lines, in the article's own notation, and no retained line is substituted or re-flowed.  Retained uncounted as the source-local context the proof needs: lines 1647--1649; lines 1652--1654.  Nothing was omitted from any retained span, so the package carries no omission record.  No retained line contains a citation, so the wrapper carries no bibliography and no bibliography entry.  No retained line contains a figure environment, an image-inclusion command or drawing code, so no image is delivered and the package declares no asset dependency.  Editorial formatting only, in addition: an AMS \texttt{amsart} preamble that restates the article's own declarations and macros used by the retained lines, namely the environments \texttt{theorem}, \texttt{lemma} on separate counters, together with the standard packages those lines need.  The retained environments print in this document as Theorem 1, Lemma 1 in order of appearance, whereas the article prints the same items with the numbering of its own full document.  The complete native source of the article as distributed in the arXiv e-print for this exact version is bundled unchanged as \texttt{sources/206162/source0.tex}.
\begin{theorem}\label{AglerdectoLeech}
Let $(k, \ell)$ be a pair of reproducing kernels such that the kernel functions $\{k_z\}$ are linearly independent. If $(k, \ell)$ admits a weak Agler decomposition, then it is a complete pseudo-Leech pair.
\end{theorem}

\begin{lemma}\label{noholmontel}
Suppose $\mathcal{H}_k$ be a separable RKHS over a set $X$ and  $\mathcal{E}, \mathcal{F}$  are separable Hilbert spaces. Given a sequence of contractive multipliers $\{\Phi_n\}\subset\textup{Mult}(\mathcal{H}_k\otimes\mathcal{E}, \mathcal{H}_k\otimes\mathcal{F})$, there exists a subsequence $\{m_n\}$ and a contractive multiplier $\Phi\in \textup{Mult}(\mathcal{H}_k\otimes\mathcal{E}, \mathcal{H}_k\otimes\mathcal{F})$ such that $(M_{\Phi_{m_n}})_n$ converges to $M_{\Phi}$ in the WOT topology. In particular, for all $x\in X$ the sequence $(\Phi_{m_n}(x))_n$ converges to $\Phi(x)$ in the WOT topology of $\mathcal{B}(\mathcal{E}, \mathcal{F}).$ 
\end{lemma}

\begin{theorem} \label{realization.}
Let $\ell$ be a reproducing kernel on a set $X$ and fix Hilbert spaces $\mathcal{F}, \mathcal{G}$ and 
a contractive multiplier $\opphi{\Phi}{X}{\mathcal{H}_{\ell}\otimes\mathcal{F}}{ \mathcal{H}_{\ell}\otimes\mathcal{G}}$, {with $\|\Phi(z)\|<1$ for all $z\in X$}. Suppose $F=\{z_1,z_2,\dots,z_n\}\subseteq X$
is finite. If   there exist Hilbert spaces $  \mathcal{K}, \mathcal{L}, \mathcal{J}$, a function $Q:X\to\mathcal{B}(\mathcal{G},\mathbb{C})$ and partially defined data $V:F\to\mathcal{B}(\mathcal{K}, \mathcal{J}), Y:F\to\mathcal{B}(\mathcal{L}, \mathcal{J})$  such that 
\[
  h(z,w):=Q(z)\big[I-\Phi(z)\Phi(w)^*\big]Q(w)^*
\]
 is non-vanishing on $X\times X$ and 
\begin{equation*} %\label{almostell}
(V(z)V(w)^*-Y(z)Y(w)^*)h^{-1}(z, w)\succeq 0, \hspace*{0.4 cm} z, w\in F,
\end{equation*}
then there exists a contractive multiplier $\opphi{\Psi}{X}{\mathcal{H}_{\ell}\otimes\mathcal{L}}{\mathcal{H}_{\ell}\otimes\mathcal{K}}$ such that $Y(z)=V(z)\Psi(z)$ for all $z\in F.$
\end{theorem}

\begin{proof}[Proof of Theorem \ref{AglerdectoLeech}]
Assume $(V_iV_j^*-Y_iY_j^*)k(z_i, z_j)\succeq 0$ over a finite set $F=\{z_1, \dots, z_n\}\subset X,$
where $V_i\in\mathcal{B}(\mathcal{K},\mathcal{J})$
and $Y_i\in\mathcal{B}(\mathcal{L},\mathcal{J})$ for all $i$ and $\mathcal{K}, \mathcal{L}, \mathcal{J}$ are (finite-dimensional) Hilbert spaces. Since $(k, \ell)$ admits a weak Agler decomposition, there exists a sequence $\{i_m\}\subset\mathbb{N}$, functions $Q_m: X\to \mathcal{B}(\mathbb{C}^{i_m}, \CC)$ and contractive  multipliers $\opphi{\Phi_m}{X}{\mathcal{H}_{\ell}\otimes\mathbb{C}^{i_m}}{\mathcal{H}_{\ell}\otimes\mathbb{C}^{i_m}}$ such that
\begin{equation*}  \ 
    h_m:=Q_m\big[I-\Phi_m\Phi_m^* \big]Q_m^*\xrightarrow{\text{\hspace{0.1 cm} \hspace{0.1 cm}}}  1/k 
\end{equation*}
 pointwise on $X\times X$. Without loss of generality, assume $\|\Phi_m(z)\|<1$ for all $m$ and $z.$ Let $\{\epsilon_q\}$ denote a decreasing null sequence. Since the kernel functions $\{k_z\}$ are linearly independent, the matrix $[k(z_i, z_j)]_{1\le i, j\le n}$ is positive-definite, and thus the same holds for 
\[(\epsilon_q +V_iV_j^*-Y_iY_j^*)k(z_i, z_j), \]
for any $q.$ Since, for any $q\ge 1,$ 
\[\big[(\epsilon_q+V_iV_j^*-Y_iY_j^*)h^{-1}_{m}(z_i, z_j)\big]_{i, j}\xrightarrow[\text{}]{\text{$m$}} \big[(\epsilon_q+V_iV_j^*-Y_iY_j^*)k(z_i,z_j) \big]_{i, j}, \]
after passing to a subsequence of $\{h_m\}$ and re-indexing, we assume 
\begin{equation*} 
\big[(\epsilon_m+V_iV_j^*-Y_iY_j^*)h^{-1}_{m}(z_i, z_j)\big]_{i, j}\succeq 0,
\end{equation*}
for all $m.$ In particular, for any $m\ge 1,$ there exists an integer $p\ge 1$ and a function $R_m:\{z_1, \dots, z_n\}\to\mathcal{B}(\mathbb{C}^{p},\CC)$ such that
\begin{align*}
&\epsilon_m +V_iV_j^*-Y_iY_j^* \notag  \\ 
=\ & Q_m(z_i)\big[I-\Phi_m(z_i)\Phi_m(z_j)^* \big]Q_m(z_j)^*\, R_m(z_i)R_m(z_j)^*. %\label{actualAgler}
\end{align*}
 {Since $\|\Phi_m(z)\|<1$ pointwise, by} Theorem \ref{realization.}  there exists a contractive multiplier $\opphi{\Psi_m}{X}{\mathcal{H}_{\ell}\otimes \mathcal{L}}{ \mathcal{H}_{\ell}\otimes \mathcal{K}\oplus\mathbb{C}}$ such that $Y_i=\begin{bmatrix}
   V_i &   \sqrt{\epsilon_m}
 \end{bmatrix}\Psi_m(z_i)$ for all $i.$ By Lemma \ref{noholmontel}, we obtain a subsequence $\{j_n\}$ and a contractive multiplier $\opphi{\Psi}{X}{\mathcal{H}_{\ell}\otimes \mathcal{L}}{\mathcal{H}_{\ell}\otimes \mathcal{K}\oplus\mathbb{C}}$ such that $\Psi_{j_n}(z)\to\Psi(z)$ WOT for all $z\in X.$ For all  $u\in\mathcal{K},$ and  $v\in\mathcal{L}$ and integers $n, i\ge 1$,
 \begin{align*}
\langle \left( Y_i- \begin{bmatrix}
   V_i &  \mathbf{0}
   \end{bmatrix}\Psi(z_i)\right ) u, v \rangle=&\langle  \left (\begin{bmatrix}
   V_i &   \sqrt{\epsilon_{j_n}}
   \end{bmatrix}\Psi_{j_n}(z_i) \right ) u, v\rangle-\langle  \begin{bmatrix}
   V_i &  \mathbf{0}
   \end{bmatrix}\Psi(z_i), v \rangle \\ 
   =&\big\langle  \big(\begin{bmatrix}
   V_i &   \sqrt{\epsilon_{j_n}}
   \end{bmatrix}-\begin{bmatrix}
   V_i &  \mathbf{0}
   \end{bmatrix}\big)\Psi_{j_n}(z_i) u, v\big\rangle \\
   & + \big\langle  \begin{bmatrix}
   V_i &  \mathbf{0}
   \end{bmatrix}\big(\Psi_{j_n}(z_i)-\Psi(z_i)\big)u, v\big\rangle.
 \end{align*}
 Since $\Psi_{j_n}(z_i)\to\Psi(z_i)$ WOT and also $\|\Psi_{m}(z)\|\le 1$ for all $m$ and $z,$ taking $n\to\infty$ yields $\langle  ( Y_i- \begin{bmatrix}
   V_i &  \mathbf{0}
   \end{bmatrix}\Psi(z_i) )u, v \rangle=0,$ for all $i$ and $u,v.$
  Thus, $Y_i=\begin{bmatrix}
   V_i &  \mathbf{0}
   \end{bmatrix}\Psi(z_i)$, for all $i$. Finally,
it is readily verified that 
 \[\Psi(z)=\begin{bmatrix}
     \Psi_1(z) \\ 
     \Psi_2(z)
 \end{bmatrix}, \]
 with $\Psi_1(z)\in\mathcal{B}(\mathcal{L}, \mathcal{K}), \Psi_2(z)\in\mathcal{B}(\mathcal{L}, \mathbb{C})$ and therefore $\Psi_1$ is a contractive multiplier  such that $Y_i=V_i\Psi_1(z_i),$ for all $i.$ Thus, $(k, \ell)$ is complete pseudo-Leech and the proof is complete. 
\end{proof}

\end{document}
